\section*{Chapitre \ref{ch:s1:residual-and-principal-component-stat-arb} --- Arbitrage statistique résiduel et en composantes principales}

\begin{solution}{exo:s1:residual-and-principal-component-stat-arb:1}
$\kappa = -252 \ln 0.9 = 26.6$ par an~; le temps de retour à la moyenne est $252/\kappa = 9.5$ jours.
\end{solution}

\begin{solution}{exo:s1:residual-and-principal-component-stat-arb:2}
$m = 0.002/0.1 = 0.02$~; $\sigma_{\mathrm{eq}} = \sqrt{10^{-4}/(1 - 0.81)} = 2.29~\%$~; $s = -0.02/0.0229 = -0.87$. Non~: $|s| < 1.25$.
\end{solution}

\begin{solution}{exo:s1:residual-and-principal-component-stat-arb:3}
$\kappa = -252 \ln b > 252/30$ équivaut à $\ln b < -1/30$, soit $b < e^{-1/30} = 0.9672$.
\end{solution}

\begin{solution}{exo:s1:residual-and-principal-component-stat-arb:4}
Le poids échangé vaut le double du taux de rotation~: $2 \times 0.37 \times 0.0005 \times 252 = 9.3~\%$ par an. Le 12.2~\% du chapitre ajoute aussi les couvertures, dont les expositions sectorielles
changent à mesure que les positions s'ouvrent et se ferment.
\end{solution}

\begin{solution}{exo:s1:residual-and-principal-component-stat-arb:5}
Les résidus d'une régression avec ordonnée à l'origine ont une somme nulle, de sorte que le résidu cumulé part et revient près de zéro. Un AR(1) ajusté sur une telle trajectoire ancrée de
soixante points a un coefficient biaisé nettement inférieur à un, de sorte que presque toutes les trajectoires semblent revenir à la moyenne en moins de trente jours, y compris une marche aléatoire.
Le filtre teste l'estimation, et l'estimation est biaisée.
\end{solution}

\begin{solution}{exo:s1:residual-and-principal-component-stat-arb:6}
L'indice sectoriel retire le marché et l'industrie de l'action, mais pas les styles~; quinze \omterm{def:s1:residual-and-principal-component-stat-arb:eigenportfolio}{portefeuilles propres} couvrent le marché, les industries et les styles, de sorte que
les résidus de l'ACP et les positions couvertes portent moins de risque commun~: 7.5~\% contre 12.1~\% de volatilité, avec à peu près le même nombre de positions.
\end{solution}

\begin{solution}{exo:s1:residual-and-principal-component-stat-arb:7}
1.22 pour une part de 0.15 et $-0.06$ pour une part nulle. La seconde mesure la méthode sur un marché sans niveau de retour à la moyenne~: son ratio de Sharpe brut de 1.30 est le
retournement à un jour injecté qui fuit dans le \omterm{def:s1:residual-and-principal-component-stat-arb:sscore}{s-score}, et des coûts de 9.8~\% par an l'annulent. C'est le plancher de la méthode.
\end{solution}

\begin{solution}{exo:s1:residual-and-principal-component-stat-arb:8}
Le ratio de Sharpe de 1.44 est une moyenne sur 1997--2007, nettement plus forte avant 2003 et de 0.9 ensuite, et les seuils ont été choisis sur une partie du même
échantillon. Les paramètres publiés ne sont pas le risque~; ce qui compte, c'est que les résidus d'aujourd'hui reviennent ou non à la moyenne. Il faut d'abord l'appliquer à des données récentes, avec les coûts de la société, et mesurer directement la
part de variance résiduelle qui revient à la moyenne.
\end{solution}

\begin{solution}{pb:s1:residual-and-principal-component-stat-arb:1}
\begin{enumerate}
\item Régressions temporelles des 60 derniers rendements de chaque action sur son fonds sectoriel («~ETF~»), ou sur les rendements des principaux \omterm{def:s1:residual-and-principal-component-stat-arb:eigenportfolio}{portefeuilles propres} («~ACP~»).
\item $Q^{(j)}_i = v^{(j)}_i/\sigma_i$ à partir des vecteurs propres de la matrice de corrélation~; 15, ou autant qu'il en faut pour expliquer 55~\% de la variance.
\item $\kappa = -252 \ln b$, $m = a/(1-b)$, $\sigma_{\mathrm{eq}} = \sqrt{\operatorname{Var}\zeta/(1-b^2)}$ à partir de $X_{n+1} = a + bX_n + \zeta$.
\item $s = (X - m)/\sigma_{\mathrm{eq}} = -m/\sigma_{\mathrm{eq}}$~; ouvrir à $\pm 1.25$, clôturer les positions longues à $-0.50$ et les positions courtes à $0.75$.
\item Un niveau de retour à la moyenne dans le rendement spécifique de chaque action, occupant une part de sa variance, avec des temps de retour à la moyenne d'environ vingt jours~; le \omterm{def:s1:residual-and-principal-component-stat-arb:sscore}{s-score} est corrélé
à $-0.128$ avec la force de rappel attendue du niveau, pour une part de 0.3.
\item Secteur~: 0.07, 0.93, 2.15. ACP~: $-0.06$, 1.22, 3.00.
\item Les \omterm{def:s1:residual-and-principal-component-stat-arb:eigenportfolio}{portefeuilles propres} retirent aussi les facteurs de style~; la volatilité du portefeuille ACP est de 7.5~\% contre 12.1~\%.
\item Rien après coûts (0.07 et $-0.06$)~; son ratio de Sharpe brut de 1.0--1.3 vient du retournement d'un jour injecté.
\item 99.3~\% avec filtre, 99.1~\% sans~: l'AR(1) sur une trajectoire ancrée de 60 points est biaisé vers un retour rapide à la moyenne.
\item Ils rejettent davantage de titres et abaissent le ratio de Sharpe net~: 2.03, 1.81 et 0.74 pour des temps inférieurs à 15, 10 et 5 jours.
\item À peine~: 2.20 à 2.30 entre des seuils de 0.75 et 1.5, 1.87 à 2.0.
\item Il abaisse le ratio de Sharpe net à 1.78~: le volume synthétique augmente avec le choc spécifique du jour, de sorte que le temps de négociation raccourcit les mouvements qui reviennent.
\item \textbf{Résultat clé.} Avec une part de retour à la moyenne de 0.3, le portefeuille sectoriel affiche un ratio net de 2.15 avec le filtre de vitesse de l'article et de 2.16 sans (ACP~: 3.00 et
3.04)~; entre des seuils d'entrée de 0.75 et 1.5, le ratio net reste dans l'intervalle 2.20--2.30, et à 2.0 il vaut 1.87.
\item Le premier vecteur propre est presque fixe (cosinus de 0.999 entre estimations mensuelles)~; le sous-espace à quinze dimensions se recoupe à 0.89 en moyenne, et sa direction la moins
stable ne se recoupe qu'à 0.63.
\item Ratios de Sharpe de 0.9 pour l'ACP sur 2003--2007, contre 1.44 sur 1997--2007~; une dégradation semblable pour la version ETF après 2002.
\item Les portefeuilles de ce type ont été débouclés ensemble, conformément au récit de Khandani et Lo.
\item Des résidus dont la composante de retour à la moyenne est assez grande pour couvrir les coûts de négociation d'un tiers du portefeuille par jour.
\item Mesurer l'autocorrélation des rendements résiduels à des horizons de quelques jours à quelques semaines, sur des données que les ajustements n'ont pas utilisées, et la corrélation hors échantillon des \omterm{def:s1:residual-and-principal-component-stat-arb:sscore}{s-scores}
avec les rendements résiduels des semaines suivantes.
\item Le modèle de coûts, puis la fenêtre d'estimation~: le seuil d'entrée et le filtre de vitesse comptent à peine.
\item Négocier des actions pour les ramener vers le niveau que leurs facteurs indiquent, lorsqu'elles s'en sont suffisamment écartées pour couvrir le coût de l'aller-retour.
\end{enumerate}
\end{solution}

\begin{solution}{iq:s1:residual-and-principal-component-stat-arb:1}
La distance du résidu cumulé d'une action à son niveau de retour à la moyenne, en écarts-types d'équilibre, issue d'un ajustement OU sur une fenêtre glissante. Vendre
(et acheter la couverture) quand elle est élevée, acheter quand elle est basse, clôturer quand elle revient vers zéro~; dimensionner à parts égales ou par un optimiseur.
\end{solution}

\begin{solution}{iq:s1:residual-and-principal-component-stat-arb:2}
L'ACP ajuste la structure commune effective du marché, styles compris, et n'a pas besoin de données de fonds, mais ses facteurs pivotent et sont difficiles à interpréter ou à couvrir à bon marché.
Les ETF sont négociables, stables et peu coûteux à couvrir, mais laissent des expositions aux styles dans les résidus et dépendent de la composition des fonds.
\end{solution}

\begin{solution}{iq:s1:residual-and-principal-component-stat-arb:3}
Pas grand-chose. La pente de l'AR(1) est biaisée à la baisse sur petits échantillons, davantage sur une trajectoire ancrée par l'ordonnée à l'origine de la régression, et son erreur type est grande~: un
temps de retour à la moyenne de 6 jours ($b = 0.85$) est compatible avec un retour bien plus lent, voire aucun. Il faut vérifier avec une estimation corrigée du biais, une fenêtre plus longue ou le
bilan hors échantillon de l'ajustement.
\end{solution}

\begin{solution}{iq:s1:residual-and-principal-component-stat-arb:4}
Expositions aux facteurs restées dans le portefeuille (styles absents de la couverture), un débouclage en situation d'encombrement (crowding) (d'autres portefeuilles aux mêmes positions, comme en août 2007), des nouvelles touchant quelques grandes
valeurs, et les couvertures~: les facteurs servant à couvrir ont-ils changé~? Décider ensuite s'il faut réduire le risque ou fournir de la liquidité.
\end{solution}

\begin{solution}{iq:s1:residual-and-principal-component-stat-arb:5}
Vectoriser~: une seule résolution par moindres carrés avec toutes les actions comme seconds membres (la même matrice de facteurs), les sommes cumulées et les moments de l'AR(1) comme opérations sur les colonnes~;
mettre à jour la matrice de corrélation et les vecteurs propres chaque semaine ou chaque mois, et non chaque jour. Cela représente quelques produits matriciels par jour.
\end{solution}

\begin{solution}{iq:s1:residual-and-principal-component-stat-arb:6}
L'AR(1) $X_{n+1} = a + bX_n + \zeta$ a pour variance stationnaire $\operatorname{Var}\zeta/(1 - b^2)$, le carré de $\sigma_{\mathrm{eq}}$. Avec une ordonnée à l'origine dans
la régression, les résidus ont une somme nulle, de sorte que $X_{60} = \sum_{n=1}^{60} \epsilon_n = 0$ et $s = (0 - m)/\sigma_{\mathrm{eq}}$.
\end{solution}
